\begin{document}
\docTitle{Tailscale setup}
\phantomsection
\label{sec:tailscale-setup}
Tailscale is installed differently depending if we are running on a machine with admin rights or without.
In case of the latter, as on the MacBook, we run a userspace \code{tailscaled}
daemon with the following flags:
\begin{DocCode}
    --tun=userspace-networking
    --state=$HOME/.tailscale/tailscaled.state
    --socket=$HOME/.tailscale/tailscaled.sock
    --socks5-server=127.0.0.1:11055
    --outbound-http-proxy-listen=127.0.0.1:11055
\end{DocCode}
To check if it is running you can run
\begin{DocCode}
    tailscale --socket "$HOME/.tailscale/tailscaled.sock" status
\end{DocCode}
Automatic starting is handled by the LaunchAgent:
\begin{DocCode}
    ~/Library/LaunchAgents/io.tailscale.tailscaled.userspace.plist
\end{DocCode}

Note that commands like \code{curl} will not use this userspace network path
out of the box. Use the local proxy, for example:
\begin{DocCode}
    curl --proxy socks5h://127.0.0.1:11055 https://example.com
\end{DocCode}
\end{document}
